\section{Print the Shortest Path}
\label{sec:graph-print-shortest-path}

\problemstatement{Given a weighted undirected graph with non-negative
weights, find not just the shortest distance from node $1$ to node $n$,
but the actual \textbf{sequence of nodes} on one shortest path.}

\begin{patternbox}
\emph{``Print/return the path,'' not just its length}, means: run
Dijkstra's as usual, but additionally remember, for each node, the
\textbf{predecessor} through which its best distance was achieved.
Following those predecessor links backward from the destination
reconstructs the path.
\end{patternbox}

\subsection*{Recording where each improvement came from}

Every time relaxation improves $\textit{dist}[v]$ using edge $u\to v$,
node $u$ is the current best ``parent'' of $v$ on the shortest path.
Store $\textit{parent}[v]=u$ at that moment. When the algorithm
finishes, $\textit{parent}$ encodes a shortest-path tree rooted at the
source. Starting from the destination and repeatedly moving to the
parent walks the path in reverse; reversing that list gives the path
from source to destination.

\begin{ideabox}[A Parent Pointer Turns Distances into Paths]
Any shortest-\emph{distance} algorithm becomes a shortest-\emph{path}
algorithm by adding one array. The parent is updated in lockstep with
the distance -- exactly when, and only when, a strictly better distance
is found -- so it always reflects the edge on the current best path.
\end{ideabox}

\subsection*{Algorithm}

\begin{enumerate}
  \item Run Dijkstra's from the source. Initialize
        $\textit{parent}[i]=i$ for all $i$.
  \item On each successful relaxation of edge $u\to v$, set
        $\textit{parent}[v]=u$.
  \item If $\textit{dist}[\textit{dest}]=\infty$, no path exists.
        Otherwise start at \textit{dest}, follow \textit{parent} until
        reaching the source, collecting nodes, then reverse.
\end{enumerate}

\subsection*{C++ implementation}

\begin{lstlisting}
#include <bits/stdc++.h>
using namespace std;

vector<int> shortestPath(int n, vector<vector<pair<int,int>>>& adj, int src, int dest) {
    vector<int> dist(n + 1, INT_MAX), parent(n + 1);
    for (int i = 1; i <= n; i++) parent[i] = i;
    dist[src] = 0;
    priority_queue<pair<int,int>, vector<pair<int,int>>, greater<>> pq;
    pq.push({0, src});

    while (!pq.empty()) {
        auto [d, node] = pq.top(); pq.pop();
        if (d > dist[node]) continue;
        for (auto& [nbr, w] : adj[node]) {
            if (dist[node] + w < dist[nbr]) {
                dist[nbr] = dist[node] + w;
                parent[nbr] = node;          // remember predecessor
                pq.push({dist[nbr], nbr});
            }
        }
    }

    if (dist[dest] == INT_MAX) return {};    // unreachable
    vector<int> path;
    for (int cur = dest; cur != src; cur = parent[cur]) path.push_back(cur);
    path.push_back(src);
    reverse(path.begin(), path.end());
    return path;
}

int main() {
    int n = 5;
    vector<vector<pair<int,int>>> adj(n + 1);   // 1-indexed
    auto addEdge = [&](int u, int v, int w) {
        adj[u].push_back({v, w});
        adj[v].push_back({u, w});
    };
    addEdge(1, 2, 2); addEdge(1, 3, 4); addEdge(2, 3, 1);
    addEdge(2, 4, 7); addEdge(3, 5, 3); addEdge(5, 4, 2);

    vector<int> path = shortestPath(n, adj, 1, 4);
    for (int node : path) cout << node << " ";
    cout << endl;                            // 1 2 3 5 4
    return 0;
}
\end{lstlisting}

\subsection*{Dry run}

\dryrun{Source $1$, destination $4$.}

\begin{itemize}
  \item Dijkstra finalizes: $\textit{dist}[1]=0$, $[2]=2$ (parent
        $1$), $[3]=3$ (parent $2$, via $2\!+\!1$), $[5]=6$ (parent
        $3$), $[4]=8$ (parent $5$, via $6\!+\!2$, better than
        $2\!+\!7=9$).
  \item Backtrack from $4$: $4\to5\to3\to2\to1$.
  \item Reverse: $1,2,3,5,4$ (total distance $8$).
\end{itemize}

\subsection*{Complexity}

\begin{complexitybox}
\textbf{Time Complexity.} $O((V+E)\log V)$ -- Dijkstra's cost; the
backtrack is $O(V)$. \textbf{Space Complexity.} $O(V)$ for the distance
and parent arrays.
\end{complexitybox}

\begin{takeawaybox}
To turn any relaxation-based shortest-distance algorithm into one that
prints the path, record a parent pointer at each successful relaxation
and walk it backward from the destination. The same trick works for
BFS, Dijkstra's, and Bellman-Ford alike.
\end{takeawaybox}
